\writeSectionFrame{\presentationTitle}{Fazit}
\realworld{Comparable}{Vergleiche bei Sortieralgorithmen}
\note{
	Um das Vergleichen z.B. in einem Sortieralgorithmus zu ermöglichen, implementieren wir die Methode compareTo() von Comparable.  Der restliche Algorithmus läuft dann immer gleich ab wie in Comparable definiert. Die spezifische Sortierlogik wird dann an die Klassen delegiert, die das Interface implementieren. Die allgemeine Ablaufsteuerung, wie diese Vergleiche dann zu einer Sortierung führen, ist in den Java-Collections-Klassen vorgegeben. Dies trennt die Verantwortlichkeiten und ermöglicht es, die allgemeine Logik der Sortierung von den spezifischen Vergleichsdetails zu trennen.
}

\begin{frame}[fragile]{Comparable-Beispiel}
	\begin{lstlisting}
public class Person implements Comparable<Person> {
    private String name;
    private int age;

    public Person(String name, int age) {
        this.name = name;
        this.age = age;
    }

    @Override
    public int compareTo(Person other) {
        return Integer.compare(this.age, other.age);
    }
}	    
	\end{lstlisting}
\end{frame}

\realworld{JavaFX}{Szenegraphen}
\note{
	JavaFX basiert bekanntermaßen auf einem Szenegraphen. Und innerhalb dieses Szenegraphen wird Schablone mehrmals verwendet. Diese Methode wird in benutzerdefinierten Pane-Unterklassen überschrieben, um die Anordnung der Kinder-Elemente innerhalb eines Containers zu bestimmen. Die allgemeine Logik des Layouts und Renderings wird durch die JavaFX-Engine gesteuert, aber spezifische Layout-Logiken werden durch das Überschreiben von layoutChildren() definiert. Bei einem benutzerdefinierten Layout würde man genau diese Methode überschreiben und sonst den von JavaFX vorgegebenen Algorithmus verwenden. In JavaFX-Tabellenzellen wird das Schablonenmuster durch die Methode updateItem(T item, boolean empty) implementiert. Diese Methode wird aufgerufen, um den Inhalt einer Zelle basierend auf dem zugeordneten Datenobjekt zu aktualisieren. Die allgemeine Ablaufsteuerung wird von JavaFX verwaltet, während die spezifische Darstellung der Zelle in updateItem() definiert wird. Und auch hier würde man bei einer benutzerdefinierten Zeilen-Darstellung eingreifen.
}

\realworld{JUnit}{Vor und nach Testfällen}
\note{
	JUnit ermöglicht es vor und nach einem Testfall bzw. allen Tests Aktionen durchzuführen. Hier wird die Schablone mit Hooks verwendet. Aber auch die allgemeine Ablaufsteuerung verwendet Schablone. Entwickler implementieren spezifische Testmethoden wie testMethodA(), während die allgemeine Ablaufsteuerung, wie Tests ausgeführt und Ergebnisse gesammelt werden, in der Basisklasse definiert ist.
}

\realworld{Maven}{Phasen des Lebenszyklus}
\note{
	In Maven werden Schablone und Hooks mehrmals verwendet. Maven definiert eine Reihe von Build-Lebenszyklen, die jeweils eine Abfolge von Phasen (Phases) enthalten. Der allgemeine Ablauf des Builds ist in diesen Phasen festgelegt und wird von Maven verwaltet. Diese Phasen sind wie die Template Method, die den grundlegenden Ablauf definiert. Innerhalb dieser Phasen können spezifische Plugins ausgeführt werden, die das Verhalten in diesen Phasen bestimmen.
}

\realworld{Java I/O}{Allgemeine Struktur für das Lesen und Schreiben von Daten}
\note{
	Das Java I/O-Framework definiert eine allgemeine Struktur für das Lesen und Schreiben von Daten, während spezifische Implementierungen bestimmte Schritte des Prozesses anpassen. 
}

\realworld{Liferay}{Anpassungen des Web-Frameworks}
\note{
	Liferay ist ein großes Web-Portal-System, das Schablone sehr viel einsetzt. Es definiert kleine Webanwendungen namens Portlets und setzt diese zu einer großen Portalanwendung zusammen. Ein Portlet in Liferay durchläuft verschiedene Phasen, und die Grundstruktur dieses Lebenszyklus wird durch die Portlet-Schnittstellen und -Klassen definiert. Entwickler können dann spezifische Details in ihren benutzerdefinierten Portlets bereitstellen und implementieren dazu bestimmte Methoden wie doView() für die spezifischen Details der Ansicht. Liferay bietet zudem ein Hook-System und ein Extender-Framework, das es Entwicklern ermöglicht, die Funktionalität von Liferay zu erweitern, indem sie in den bestehenden Code eingreifen. Liferay verwendet Templates für die Erstellung und Verwaltung von Webinhalten. Das Schablonenmuster wird hier genutzt, um die Struktur von Content-Templates zu definieren, während spezifische Inhalte durch die benutzerdefinierten Templates bereitgestellt werden.
}

\begin{frame}{Wann ist der Einsatz des Schablonenmusters sinnvoll?}
\begin{itemize}
    \item Wiederkehrende Algorithmen mit variablen Schritten
    \item Vermeidung von Code-Duplikation
    \item Erzwingen eines festen Ablaufs
    \item Verwaltung komplexer Prozesse
    \item Ermöglichung von Erweiterungen
    \item Integration in bestehende Frameworks
\end{itemize}
\end{frame}

\note{
	 Im Grunde kann die Schablonenmethode angewendet werden bei generischen Algorithmen, 
	 deren Verhalten teilweise variiert werden kann, aber nicht unbedingt variiert werden
	 muss. 
}

\vorteil{Austauschbarkeit von Teilen von Algorithmen}
\vorteil{Wiederverwendbarkeit von Teilen von Algorithmen}
\note{
    Das Schablonenmuster ermöglicht es, die allgemeine Struktur eines Algorithmus in einer Basisklasse zu definieren, während Unterklassen spezifische Schritte anpassen oder implementieren können. Dies fördert die Wiederverwendung von Code und reduziert Redundanz.
}

\vorteil{Einheitlichkeit}
\vorteil{Klare Definition von Stellen, an denen Verhalten geändert werden kann}
\note{
    Da die allgemeine Struktur des Algorithmus in der Basisklasse definiert ist, bleibt der Ablauf konsistent, selbst wenn Unterklassen verschiedene Implementierungen für einzelne Schritte bereitstellen.
}

\vorteil{Erweiterbarkeit}
\note{
    Das Muster ermöglicht es, das Verhalten von Algorithmen zu erweitern, indem man neue Unterklassen erstellt, die spezifische Implementierungen der Methoden bereitstellen, ohne den allgemeinen Ablauf zu verändern.
}

\vorteil{Vermeidung von Duplikaten}
\note{
    Gemeinsame Schritte des Algorithmus müssen nicht in jeder Unterklasse neu implementiert werden. Dadurch wird der Code einfacher zu warten und zu erweitern.
}

\vorteil{Klarheit und Struktur}
\note{
    Das Muster trennt die allgemeine Logik von den Details der Implementierung, was den Code verständlicher und leichter zu verfolgen macht.
}

\vorteil{Open-Closed-Prinzip bei neuen konkreten Implementierungen von Algorithmen}
\note{
    Da neue Implementireungen einfach hinzugefügt werden können ohne bestehenden Code verändern zu müssen, ist das Open-Closed-Prinzip bei neuen konkreten Implementierungen von Algorithmen erfüllt. 
}

\vorteil{Single-Responsibility-Prinzip bei korrekter Implementierung erfüllt}
\note{
    Auch das Single-Responsibility-Prinzip ist erfüllt, wenn das Muster ordentlich implementiert wird. Die Basisklasse hat die Verantwortung, die allgemeine Struktur eines Algorithmus zu definieren, während die Unterklassen spezifische Implementierungen von einzelnen Schritten übernehmen. Solange diese Verantwortung klar abgegrenzt ist und die Klasse keine zusätzlichen Aufgaben übernimmt, erfüllt das Schablonenmuster das Prinzip. Das Muster trennt die allgemeine Steuerung des Algorithmus (in der Basisklasse) von der spezifischen Implementierung der einzelnen Schritte (in den Unterklassen). 
}

\nachteil{Vorgegebenes Skelett eines Algorithmus kann einige Klienten beschränken.}
\note{
    Da der Ablauf des Algorithmus in der Basisklasse festgelegt ist, kann dies die Flexibilität einschränken, wenn eine völlig andere Reihenfolge von Schritten oder eine alternative Vorgehensweise benötigt wird.
}

\nachteil{Vermehrte Anzahl von Klassen }
\note{
    Da für jede spezifische Implementierung eine neue Unterklasse erstellt werden muss, kann das Schablonenmuster zu einer größeren Anzahl von Klassen führen, was die Komplexität des Systems erhöht.
}

\nachteil{Versteckte Abhängigkeiten}
\note{
    Die Abhängigkeit zwischen der Basisklasse und den Unterklassen kann zu einer schwer nachvollziehbaren Architektur führen, insbesondere wenn der Ablauf des Algorithmus nicht klar dokumentiert ist.
}

\nachteil{Overhead durch Vererbung}
\note{
    Wenn das Schablonenmuster in einem System übermäßig eingesetzt wird, kann die extensive Verwendung von Vererbung den Code schwerer verständlich und wartbar machen. Es kann schwieriger werden, die Auswirkungen von Änderungen in der Basisklasse auf die Unterklassen zu verstehen und zu kontrollieren.
}

\nachteil{Schwierige Fehlerbehebung}
\note{
    Da die eigentliche Implementierung von Methoden in den Unterklassen erfolgt und die Ablaufsteuerung in der Basisklasse, kann es schwieriger sein, Fehler zu lokalisieren und zu beheben, insbesondere in großen und komplexen Hierarchien.
}

\nachteil{Schwere Wartung bei vielen Einhakmöglichkeiten}
\note{
    Vor allem der Hooking-Mechanismus sollte mit Bedacht verwendet werden. Es ergibt keinen Sinn alle möglichen Haken anzubieten, denn damit steigt der Wartungsaufwand. Man sollte sich wenige sinnvolle Einhakmöglichkeiten an sinnvollen Stellen überlegen.
}

\nachteil{Liskov-Substitutionsprinzip könnte durch Verhaltensänderungen gebrochen werden}
\note{
    Das Liskov-Substitutionsprinzip könnte gebrochen werden, wenn Unterklassen das Verhalten der Basisklasse zu sehr verändern. 
}

%\begin{frame}{Verwandte oder ergänzende Muster}
%	\begin{itemize}
% 		\item Fabrikmethode kann als Schablonenmethode implementiert werden
% 		\item Verwandtschaft zu Strategie
% 	\end{itemize}
%\end{frame}

%\note{
%	Die Schablonenmethode basiert auf Vererbung und ermöglicht das Austauschen von Teilen eines
%	Algorithmus, während Strategie auf Komposition basiert und ermöglicht das Austauschen von
%	ganzen Algorithmen. Die Schablonenmethode arbeitet auf Klassenleven und ist somit statisch,
%	während Strategie auf Objektebene arbeitet und einen Austausch während der Laufzeit
%	ermöglicht.
%}